\documentclass[10pt,letterpaper,sans]{moderncv}
\moderncvstyle{banking}
\moderncvcolor{blue}

\usepackage[scale=0.81]{geometry}
\nopagenumbers{}
\setlength{\hintscolumnwidth}{2.25cm}

\name{Krishnaa}{Vadivel}
\title{Computational Physics | Distributed Computing | Machine Learning}
\email{srikrishnaa.careers@gmail.com}
\phone[mobile]{516-853-5663}
\social[linkedin]{sri-krishnaa-ece-phd}
\extrainfo{\href{https://krishnaa423.github.io/personal-website/}{Website} \textbar{} \href{https://github.com/krishnaa423}{GitHub: krishnaa423} \textbar{} \href{https://leetcode.com/u/krishnaa42342}{LeetCode: krishnaa42342} \textbar{} \href{https://hub.docker.com/u/krishnaa42342}{Docker Hub: krishnaa42342}}

\begin{document}
\makecvtitle

\section{Summary}
\cvitem{}{Computational physics researcher building distributed scientific software, machine-learning accelerators, and first-principles materials workflows. Experience spans Python/C++, CUDA/ROCm, MPI/OpenMP, Linux-based HPC environments, custom Docker containers for reproducibility, and scientific ML for excitonic materials problems.}

\section{Research Experience}
\cventry{2021--present}{PhD Research Assistant}{Yale University, Electrical Engineering}{}{}{\begin{itemize}%
\item Developed first-principles and many-body theory workflows for self-trapped excitons and excitonic polarons in real materials, connecting DFT, DFPT, GW/BSE-style methods, and reduced effective models.
\item Built distributed simulation software in Python and C++ for leadership-class HPC systems, using MPI, OpenMP, CUDA, ROCm, PETSc, and SLEPc on Frontier, Perlmutter, and other national-lab clusters.
\item Formulated coherent-state and path-integral viewpoints for exciton-phonon dynamics and used PyTorch, PyTorch Geometric, and equivariant graph neural networks to accelerate self-trapped-exciton calculations and related materials modeling workflows.
\item Packaged reproducible scientific software in custom Docker containers and debugged distributed Linux applications with Bash and command-line tooling across HPC environments.
\item Co-authored a 2025 \textit{Journal of the American Chemical Society} paper on ultrafast self-trapped exciton formation in lead-free double halide perovskites and presented related APS research in 2024, 2025, and 2026.
\end{itemize}}

\section{Industry Experience}
\cventry{2019--2021}{Software and Embedded Systems Engineer}{Probe Technologies Inc. / Weatherford Wireline Services}{}{}{\begin{itemize}%
\item Developed CUDA-based software for high-volume acoustic waveform processing, visualization, and field-tool diagnostics in geological well logging systems.
\item Built GUI applications with C\#, ASP.NET server-side services, and Microsoft SQL Server workflows to automate remote calibration of temperature and pressure sensors in real time, improving production efficiency.
\item Developed low-level debugging tools and engineering software for deployed instrumentation and server-side programs in hardware-adjacent production environments.
\end{itemize}}
\cventry{2018--2019}{Photonics Technical Writer}{FindLight}{}{}{\begin{itemize}%
\item Wrote clear technical content on lasers, optics, and photonics devices for engineering audiences and product-focused readers.
\end{itemize}}

\section{Selected Projects}
\cvitem{Exciton-Phonon HPC Research}{Developed distributed first-principles calculations for self-trapped excitons, many-body theory workflows, and scalable materials simulations.}
\cvitem{Equivariant Scientific ML}{Built PyTorch and PyTorch Geometric workflows with equivariant graph neural networks to accelerate molecular and excitonic modeling tasks.}
\cvitem{Meep Ring Resonator FDTD}{Designed a nanophotonic ring-resonator band-pass filter in telecom bands using Meep-based FDTD simulation workflows.}
\cvitem{FEniCS-Based FEA}{Built FEniCS-based finite-element workflows for frequency analysis of steel pipes in well environments to support physics-driven engineering studies.}

\section{Education}
\cventry{2021--present}{PhD, Electrical Engineering}{Yale University}{}{}{}
\cventry{2023}{MPhil, Electrical Engineering}{Yale University}{}{}{}
\cventry{2023}{MS, Electrical Engineering}{Yale University}{}{}{}
\cventry{2017--2018}{MEng, Electrical and Computer Engineering}{Cornell University}{}{}{}
\cventry{2013--2017}{BS, Electrical and Computer Engineering}{Cornell University}{}{}{}

\section{Technical Skills}
\cvitem{Programming}{Python, C, C++, CUDA, JavaScript, Bash, SQL, C\#}
\cvitem{Scientific Computing}{MPI, OpenMP, PETSc, SLEPc, NumPy, SciPy, SymPy, matplotlib, FEniCS}
\cvitem{Machine Learning}{PyTorch, PyTorch Geometric, scikit-learn, equivariant graph neural networks}
\cvitem{Platforms}{Linux, macOS, Windows, Docker, Docker Compose, Kubernetes, gcloud, Frontier, Perlmutter}
\cvitem{Domains}{First-principles materials modeling, many-body theory, distributed computing, computational physics, numerical PDEs}

\end{document}
